\documentclass[12pt,]{article}
\usepackage{lmodern}
\usepackage{amssymb,amsmath}
\usepackage{ifxetex,ifluatex}
\usepackage{fixltx2e} % provides \textsubscript
\ifnum 0\ifxetex 1\fi\ifluatex 1\fi=0 % if pdftex
  \usepackage[T1]{fontenc}
  \usepackage[utf8]{inputenc}
\else % if luatex or xelatex
  \ifxetex
    \usepackage{mathspec}
  \else
    \usepackage{fontspec}
  \fi
  \defaultfontfeatures{Ligatures=TeX,Scale=MatchLowercase}
\fi
% use upquote if available, for straight quotes in verbatim environments
\IfFileExists{upquote.sty}{\usepackage{upquote}}{}
% use microtype if available
\IfFileExists{microtype.sty}{%
\usepackage{microtype}
\UseMicrotypeSet[protrusion]{basicmath} % disable protrusion for tt fonts
}{}
\usepackage[margin=1in]{geometry}
\usepackage{hyperref}
\hypersetup{unicode=true,
            pdftitle={Exceedingly Simple Monotone Regression (with Ties)},
            pdfauthor={Jan de Leeuw},
            pdfborder={0 0 0},
            breaklinks=true}
\urlstyle{same}  % don't use monospace font for urls
\usepackage{color}
\usepackage{fancyvrb}
\newcommand{\VerbBar}{|}
\newcommand{\VERB}{\Verb[commandchars=\\\{\}]}
\DefineVerbatimEnvironment{Highlighting}{Verbatim}{commandchars=\\\{\}}
% Add ',fontsize=\small' for more characters per line
\usepackage{framed}
\definecolor{shadecolor}{RGB}{248,248,248}
\newenvironment{Shaded}{\begin{snugshade}}{\end{snugshade}}
\newcommand{\KeywordTok}[1]{\textcolor[rgb]{0.13,0.29,0.53}{\textbf{{#1}}}}
\newcommand{\DataTypeTok}[1]{\textcolor[rgb]{0.13,0.29,0.53}{{#1}}}
\newcommand{\DecValTok}[1]{\textcolor[rgb]{0.00,0.00,0.81}{{#1}}}
\newcommand{\BaseNTok}[1]{\textcolor[rgb]{0.00,0.00,0.81}{{#1}}}
\newcommand{\FloatTok}[1]{\textcolor[rgb]{0.00,0.00,0.81}{{#1}}}
\newcommand{\ConstantTok}[1]{\textcolor[rgb]{0.00,0.00,0.00}{{#1}}}
\newcommand{\CharTok}[1]{\textcolor[rgb]{0.31,0.60,0.02}{{#1}}}
\newcommand{\SpecialCharTok}[1]{\textcolor[rgb]{0.00,0.00,0.00}{{#1}}}
\newcommand{\StringTok}[1]{\textcolor[rgb]{0.31,0.60,0.02}{{#1}}}
\newcommand{\VerbatimStringTok}[1]{\textcolor[rgb]{0.31,0.60,0.02}{{#1}}}
\newcommand{\SpecialStringTok}[1]{\textcolor[rgb]{0.31,0.60,0.02}{{#1}}}
\newcommand{\ImportTok}[1]{{#1}}
\newcommand{\CommentTok}[1]{\textcolor[rgb]{0.56,0.35,0.01}{\textit{{#1}}}}
\newcommand{\DocumentationTok}[1]{\textcolor[rgb]{0.56,0.35,0.01}{\textbf{\textit{{#1}}}}}
\newcommand{\AnnotationTok}[1]{\textcolor[rgb]{0.56,0.35,0.01}{\textbf{\textit{{#1}}}}}
\newcommand{\CommentVarTok}[1]{\textcolor[rgb]{0.56,0.35,0.01}{\textbf{\textit{{#1}}}}}
\newcommand{\OtherTok}[1]{\textcolor[rgb]{0.56,0.35,0.01}{{#1}}}
\newcommand{\FunctionTok}[1]{\textcolor[rgb]{0.00,0.00,0.00}{{#1}}}
\newcommand{\VariableTok}[1]{\textcolor[rgb]{0.00,0.00,0.00}{{#1}}}
\newcommand{\ControlFlowTok}[1]{\textcolor[rgb]{0.13,0.29,0.53}{\textbf{{#1}}}}
\newcommand{\OperatorTok}[1]{\textcolor[rgb]{0.81,0.36,0.00}{\textbf{{#1}}}}
\newcommand{\BuiltInTok}[1]{{#1}}
\newcommand{\ExtensionTok}[1]{{#1}}
\newcommand{\PreprocessorTok}[1]{\textcolor[rgb]{0.56,0.35,0.01}{\textit{{#1}}}}
\newcommand{\AttributeTok}[1]{\textcolor[rgb]{0.77,0.63,0.00}{{#1}}}
\newcommand{\RegionMarkerTok}[1]{{#1}}
\newcommand{\InformationTok}[1]{\textcolor[rgb]{0.56,0.35,0.01}{\textbf{\textit{{#1}}}}}
\newcommand{\WarningTok}[1]{\textcolor[rgb]{0.56,0.35,0.01}{\textbf{\textit{{#1}}}}}
\newcommand{\AlertTok}[1]{\textcolor[rgb]{0.94,0.16,0.16}{{#1}}}
\newcommand{\ErrorTok}[1]{\textcolor[rgb]{0.64,0.00,0.00}{\textbf{{#1}}}}
\newcommand{\NormalTok}[1]{{#1}}
\usepackage{graphicx,grffile}
\makeatletter
\def\maxwidth{\ifdim\Gin@nat@width>\linewidth\linewidth\else\Gin@nat@width\fi}
\def\maxheight{\ifdim\Gin@nat@height>\textheight\textheight\else\Gin@nat@height\fi}
\makeatother
% Scale images if necessary, so that they will not overflow the page
% margins by default, and it is still possible to overwrite the defaults
% using explicit options in \includegraphics[width, height, ...]{}
\setkeys{Gin}{width=\maxwidth,height=\maxheight,keepaspectratio}
\IfFileExists{parskip.sty}{%
\usepackage{parskip}
}{% else
\setlength{\parindent}{0pt}
\setlength{\parskip}{6pt plus 2pt minus 1pt}
}
\setlength{\emergencystretch}{3em}  % prevent overfull lines
\providecommand{\tightlist}{%
  \setlength{\itemsep}{0pt}\setlength{\parskip}{0pt}}
\setcounter{secnumdepth}{5}
% Redefines (sub)paragraphs to behave more like sections
\ifx\paragraph\undefined\else
\let\oldparagraph\paragraph
\renewcommand{\paragraph}[1]{\oldparagraph{#1}\mbox{}}
\fi
\ifx\subparagraph\undefined\else
\let\oldsubparagraph\subparagraph
\renewcommand{\subparagraph}[1]{\oldsubparagraph{#1}\mbox{}}
\fi

%%% Use protect on footnotes to avoid problems with footnotes in titles
\let\rmarkdownfootnote\footnote%
\def\footnote{\protect\rmarkdownfootnote}

%%% Change title format to be more compact
\usepackage{titling}

% Create subtitle command for use in maketitle
\newcommand{\subtitle}[1]{
  \posttitle{
    \begin{center}\large#1\end{center}
    }
}

\setlength{\droptitle}{-2em}
  \title{Exceedingly Simple Monotone Regression (with Ties)}
  \pretitle{\vspace{\droptitle}\centering\huge}
  \posttitle{\par}
  \author{Jan de Leeuw}
  \preauthor{\centering\large\emph}
  \postauthor{\par}
  \predate{\centering\large\emph}
  \postdate{\par}
  \date{Version 01, April 01, 2017}


\begin{document}
\maketitle
\begin{abstract}
A C implementation of Kruskal's up-and-down-blocks monotone regression
algorithm for use with .C() is extended to include the three classic
ways of handling ties. It is then compared with other implementations.
\end{abstract}

{
\setcounter{tocdepth}{3}
\tableofcontents
}
Note: This is a working paper which will be expanded/updated frequently.
All suggestions for improvement are welcome. The directory
\href{http://gifi.stat.ucla.edu/jbkTies}{gifi.stat.ucla.edu/jbkTies} has
a pdf version, the bib file, the complete Rmd file with the code chunks,
and the R source code.

\section{Introduction}\label{introduction}

In a recent Rpub (De Leeuw (2017a)) we presented a \texttt{.C()} version
of Kruskal's up-and-down-blocks algorithm for monotone (or isotone)
regression. In De Leeuw (2016) there is R code that extends monotone
regression by adding the three classical ways of handling ties. See
Kruskal (1964) for the primary and secondary approach, and De Leeuw
(1977) for the tertiary approach, as well as for proofs that all three
methods are indeed optimal. The R code for tie handling in De Leeuw
(2016) uses lists for blocks of ties, and includes a number of loops
over blocks. Thus there is some room for improvement by implementing
tie-handling using \texttt{.C()}.

The appendix of this follow-up paper presents \texttt{.C()} code for
monotone regression with the three approaches. We have chosen to make it
modular, in the sense that the basic operations all have a separate
routine written in C, with correspondng R wrapper. All C routines are
written using \texttt{.C()} conventions, i.e.~they pass by reference and
return a void. Internally they do not use any R-based functions and
include files, so they can be easily used in other contexts as well.
Thus extra storage need in the routines is allocated and freed
internally, without involving R.

\section{Small Example}\label{small-example}

\begin{Shaded}
\begin{Highlighting}[]
\NormalTok{x <-}\StringTok{ }\KeywordTok{c}\NormalTok{(}\FloatTok{2.1}\NormalTok{, }\FloatTok{2.1}\NormalTok{, }\FloatTok{3.5}\NormalTok{, }\FloatTok{1.9}\NormalTok{, }\FloatTok{3.5}\NormalTok{, }\FloatTok{3.5}\NormalTok{, }\FloatTok{1.9}\NormalTok{, }\FloatTok{2.1}\NormalTok{, }\FloatTok{1.9}\NormalTok{)}
\NormalTok{y <-}\StringTok{ }\KeywordTok{c}\NormalTok{(}\DecValTok{2}\NormalTok{, }\DecValTok{1}\NormalTok{, }\DecValTok{6}\NormalTok{, }\DecValTok{5}\NormalTok{, }\DecValTok{4}\NormalTok{, }\DecValTok{7}\NormalTok{, }\DecValTok{8}\NormalTok{, }\DecValTok{9}\NormalTok{, }\DecValTok{3}\NormalTok{)}
\NormalTok{w <-}\StringTok{ }\KeywordTok{c}\NormalTok{(}\DecValTok{1}\NormalTok{, }\DecValTok{1}\NormalTok{, }\DecValTok{2}\NormalTok{, }\DecValTok{2}\NormalTok{, }\DecValTok{2}\NormalTok{, }\DecValTok{2}\NormalTok{, }\DecValTok{2}\NormalTok{, }\DecValTok{1}\NormalTok{, }\DecValTok{1}\NormalTok{)}
\end{Highlighting}
\end{Shaded}

\subsection{Sorting}\label{sorting}

We use the function \texttt{mySortDouble()}, a \texttt{.C()} interface
to a C routine of the same name, to sort x and to put y and w in the
correct order. In R parlance, the routine returns \texttt{order(x)} as
well as \texttt{sort\ (x)\ =\ x{[}order\ (x){]}} ,
\texttt{y{[}order\ (x){]}} and \texttt{w{[}order\ (x){]}}. The sorting
routine is C code taken from De Leeuw (2017b). It uses the system
\texttt{qsort()} routine, which implements quicksort. Quicksort is not a
stable sorting routine, in the sense that it does not guarantee that
tieblocks will be sorted in a unique way. This, however, is not a
problem in the monotone regression context, because that form of
instability will not change the outcome of the regression routine.

\begin{verbatim}
## [1] 1.9 1.9 1.9 2.1 2.1 2.1 3.5 3.5 3.5
\end{verbatim}

\begin{verbatim}
## [1]   8   3   5   2   1   9   6   7   4
\end{verbatim}

\begin{verbatim}
## [1]   2   1   2   1   1   1   2   2   2
\end{verbatim}

\begin{verbatim}
## [1]   7   9   4   1   2   8   3   6   5
\end{verbatim}

The function \texttt{tieBlock()}, again a \texttt{.C()} interface to a C
routine of the same name, uses the sorted x to return the tie-blocks (a
vector of the same length as x, using integers to indicate block
membership).

\begin{verbatim}
## [1]   1   1   1   2   2   2   3   3   3
\end{verbatim}

We emphasize that in a non-metric multidimensional scaling context (or,
more generally, an alternating least squares context) the sorting only
has to be done only once, not in every iteration. Thus in repeated calls
of the algorithm this first part can be skipped.

\subsection{Primary Approach}\label{primary-approach}

The primary approach to ties starts with sorting y within blocks. This
is done with \texttt{sortBlocks()}, again a \texttt{.C()} interface.
Along with y we sort w and the order vector.

\subsubsection{Sort within Blocks}\label{sort-within-blocks}

\begin{verbatim}
## [1]   3   5   8   1   2   9   4   6   7
\end{verbatim}

\begin{verbatim}
## [1]   1   2   2   1   1   1   2   2   2
\end{verbatim}

\begin{verbatim}
## [1]   9   4   7   2   1   8   5   3   6
\end{verbatim}

The newly sorted y and w are entered into \texttt{jbkPava()}, the
monotone regression routine, again a \texttt{.C()} interface to C code.

\subsubsection{Monotone Regression}\label{monotone-regression}

\begin{verbatim}
## [1] 3.000 4.833 4.833 4.833 4.833 5.667 5.667 6.000 7.000
\end{verbatim}

\subsubsection{Restore Order}\label{restore-order}

Finally \texttt{mySortInteger()} ,another \texttt{.C()} interface, is
used to find the inverse permutation that puts the monotone regression
output back in the correct order.

\begin{verbatim}
## [1]   5   4   8   2   7   9   3   6   1
\end{verbatim}

\begin{verbatim}
## [1] 4.833 4.833 6.000 4.833 5.667 7.000 4.833 5.667 3.000
\end{verbatim}

\subsubsection{Residual Sum of Squares}\label{residual-sum-of-squares}

Equal to 46.6666666667.

\subsection{Secondary Approach}\label{secondary-approach}

\subsubsection{Make Blocks}\label{make-blocks}

In the secondary approach we use \texttt{makeBlocks()}, another
\texttt{.C()} interface, to create block averages and block weights.

\begin{verbatim}
## [1] 5.800 4.000 5.667
\end{verbatim}

\begin{verbatim}
## [1] 5.000 3.000 6.000
\end{verbatim}

\subsubsection{Monotone Regression}\label{monotone-regression-1}

Then we perform monotone regression on the block values using block
weights.

\begin{verbatim}
## [1] 5.125 5.125 5.667
\end{verbatim}

\subsubsection{Expand and Restore Order}\label{expand-and-restore-order}

Finally we use \texttt{mySortInteger()}, another \texttt{.C()} sorting
routine, to expand the monotone regression values in the correct order
to a vector of \texttt{length\ (x)}.

\begin{verbatim}
## [1] 5.125 5.125 5.667 5.125 5.667 5.667 5.125 5.125 5.125
\end{verbatim}

\subsubsection{Residual Sum of Squares}\label{residual-sum-of-squares-1}

Equal to 59.2604166667.

\subsection{Tertiary Approach}\label{tertiary-approach}

\subsubsection{Adjust Means}\label{adjust-means}

The tertiary approach forms blocks and does a monotone regression as in
the secondary approach. It then adjusts the block means optimally.

\begin{verbatim}
## [1] 7.325  2.325  4.325  3.125  2.125  10.125 6.000  7.000  4.000
\end{verbatim}

\subsubsection{Restore Order}\label{restore-order-1}

\begin{verbatim}
## [1] 3.125  2.125  6.000  4.325  4.000  7.000  7.325  10.125 2.325
\end{verbatim}

\subsubsection{Residual Sum of Squares}\label{residual-sum-of-squares-2}

Equal to 5.16375.

\section{Timings}\label{timings}

We compare two different implementations of monotone regression with
ties. The first is \texttt{monregR()}, which does the data-handling and
tie-handling in R, using lists. It is basically the code from De Leeuw
(2016), with two changes. In the first place the Fortran code for
monotone regression from Cran (1980) is replaced by the C code from De
Leeuw (2017a). Secondly, the R code in De Leeuw (2016) does not work
correctly in the case of unequal weights, and we have corrected that in
the current version of \texttt{monregR()}. The second function is
\texttt{monregRC()}, which does data-handling, tie-handling, and
monotone regression in C, using the functions we have illustrated
earlier in our small example.

In our comparion we use

\begin{Shaded}
\begin{Highlighting}[]
\NormalTok{x <-}\StringTok{ }\KeywordTok{sample} \NormalTok{(}\DecValTok{1}\NormalTok{:blks, n, }\DataTypeTok{replace =} \OtherTok{TRUE}\NormalTok{)}
\NormalTok{y <-}\StringTok{ }\KeywordTok{rnorm} \NormalTok{(n)}
\end{Highlighting}
\end{Shaded}

where n = 10,000. Since we sample with equal probabilities from
\texttt{1:blks} the vector x will have (at most) blks different values,
so if blks is small there will be many ties, if blks is large there will
be only a few ties. We let blks take the values
\texttt{c(2,10,100,1000,10000)}. Combined with the three options for
tie-handling this gives 15 different analyses, and each one is repeated
100 times. We report the elapsed time from \texttt{system.time()}.

Results for \texttt{monregR()} are

\begin{verbatim}
##       1       2       3      
## 2     0.3140  0.1690  0.1810 
## 10    0.3330  0.2370  0.2430 
## 100   0.9480  0.6040  0.6490 
## 1000  6.0610  3.9600  3.8570 
## 10000 36.5820 23.8970 24.3780
\end{verbatim}

Results for \texttt{monregRC()} are

\begin{verbatim}
##       1      2      3     
## 2     1.1110 1.0080 0.9850
## 10    0.9690 0.8570 0.8790
## 100   0.7640 0.6830 0.6530
## 1000  0.5070 0.4470 0.4510
## 10000 0.5060 0.4680 0.4690
\end{verbatim}

Choosing n = 10,000 is similar to what we used in De Leeuw (2017a), but
in a typical multidimensional scaling context n = 500 is perhaps more
realistic.

Now the results for \texttt{monregR()} are

\begin{verbatim}
##     1      2      3     
## 2   0.0720 0.0160 0.0200
## 10  0.0490 0.0250 0.0290
## 100 0.2500 0.0860 0.0790
## 250 0.4780 0.1850 0.1560
## 500 0.6740 0.2140 0.2210
\end{verbatim}

and those for \texttt{monregRC()} are

\begin{verbatim}
##     1      2      3     
## 2   0.0660 0.0400 0.0380
## 10  0.0350 0.0280 0.0310
## 100 0.0260 0.0210 0.0240
## 250 0.0230 0.0210 0.0240
## 500 0.0240 0.0250 0.0240
\end{verbatim}

\section{Conclusion}\label{conclusion}

Running time for \texttt{monregR()} heavily depends on the number of
blocks, because the R code has loops over blocks. For a small number of
blocks \texttt{monregR()} is even faster than \texttt{monregRC()}. If n=
10,000 then up to five times as fast for an expected block size of 5000
(two blocks). But for more blocks \texttt{monregR()} rapidly loses its
advantage. For n = 500 we find that \texttt{monregRC()} is almost always
faster, and usually much faster.

\section{Code}\label{code}

\subsection{R code}\label{r-code}

\begin{Shaded}
\begin{Highlighting}[]
\KeywordTok{dyn.load}\NormalTok{(}\StringTok{"jbkTies.so"}\NormalTok{)}

\NormalTok{mySortDouble <-}\StringTok{ }\NormalTok{function (x, y, }\DataTypeTok{w =} \KeywordTok{rep}\NormalTok{(}\DecValTok{1}\NormalTok{, }\KeywordTok{length}\NormalTok{(x))) \{}
  \NormalTok{n <-}\StringTok{ }\KeywordTok{length} \NormalTok{(x)}
  \NormalTok{xind <-}\StringTok{ }\KeywordTok{rep} \NormalTok{(}\DecValTok{0}\NormalTok{, n)}
  \NormalTok{h <-}
\StringTok{    }\KeywordTok{.C}\NormalTok{(}
      \StringTok{"mySortDouble"}\NormalTok{,}
      \DataTypeTok{x =} \KeywordTok{as.double} \NormalTok{(x),}
      \DataTypeTok{y =} \KeywordTok{as.double} \NormalTok{(y),}
      \DataTypeTok{w =} \KeywordTok{as.double} \NormalTok{(w),}
      \DataTypeTok{xind =} \KeywordTok{as.integer} \NormalTok{(xind),}
      \DataTypeTok{n =} \KeywordTok{as.integer}\NormalTok{(n)}
    \NormalTok{)}
  \KeywordTok{return} \NormalTok{(h)}
\NormalTok{\}}

\NormalTok{mySortInteger <-}\StringTok{ }\NormalTok{function (x) \{}
  \NormalTok{n <-}\StringTok{ }\KeywordTok{length} \NormalTok{(x)}
  \NormalTok{xind <-}\StringTok{ }\KeywordTok{rep} \NormalTok{(}\DecValTok{0}\NormalTok{, n)}
  \NormalTok{h <-}
\StringTok{    }\KeywordTok{.C}\NormalTok{(}
      \StringTok{"mySortInteger"}\NormalTok{,}
      \DataTypeTok{x =} \KeywordTok{as.integer} \NormalTok{(x),}
      \DataTypeTok{xind =} \KeywordTok{as.integer} \NormalTok{(xind),}
      \DataTypeTok{n =} \KeywordTok{as.integer}\NormalTok{(n)}
    \NormalTok{)}
  \KeywordTok{return} \NormalTok{(h)}
\NormalTok{\}}

\NormalTok{tieBlock <-}\StringTok{ }\NormalTok{function (x) \{}
  \NormalTok{n <-}\StringTok{ }\KeywordTok{length} \NormalTok{(x)}
  \NormalTok{h <-}
\StringTok{    }\KeywordTok{.C}\NormalTok{(}
      \StringTok{"tieBlock"}\NormalTok{,}
      \DataTypeTok{x =} \KeywordTok{as.double} \NormalTok{(x),}
      \DataTypeTok{iblks =} \KeywordTok{as.integer} \NormalTok{(}\KeywordTok{rep}\NormalTok{(}\DecValTok{0}\NormalTok{, n)),}
      \KeywordTok{as.integer} \NormalTok{(n),}
      \DataTypeTok{nblk =} \KeywordTok{as.integer} \NormalTok{(}\DecValTok{0}\NormalTok{)}
    \NormalTok{)}
  \KeywordTok{return} \NormalTok{(h)}
\NormalTok{\}}

\NormalTok{makeBlocks <-}\StringTok{ }\NormalTok{function  (x, }\DataTypeTok{w =} \KeywordTok{rep} \NormalTok{(}\DecValTok{1}\NormalTok{, }\KeywordTok{length}\NormalTok{(x)), iblks) \{}
  \NormalTok{n <-}\StringTok{ }\KeywordTok{length} \NormalTok{(x)}
  \NormalTok{nblk <-}\StringTok{ }\KeywordTok{max} \NormalTok{(iblks)}
  \NormalTok{xblks <-}\StringTok{ }\KeywordTok{rep} \NormalTok{(}\DecValTok{0}\NormalTok{, nblk)}
  \NormalTok{wblks <-}\StringTok{ }\KeywordTok{rep} \NormalTok{(}\DecValTok{0}\NormalTok{, nblk)}
  \NormalTok{h <-}
\StringTok{    }\KeywordTok{.C}\NormalTok{(}
      \StringTok{"makeBlocks"}\NormalTok{,}
      \DataTypeTok{x =} \KeywordTok{as.double} \NormalTok{(x),}
      \DataTypeTok{w =} \KeywordTok{as.double} \NormalTok{(w),}
      \DataTypeTok{xblks =} \KeywordTok{as.double} \NormalTok{(xblks),}
      \DataTypeTok{wblks =} \KeywordTok{as.double} \NormalTok{(wblks),}
      \DataTypeTok{iblks =} \KeywordTok{as.integer} \NormalTok{(iblks),}
      \DataTypeTok{n =} \KeywordTok{as.integer}\NormalTok{(n),}
      \DataTypeTok{nblk =} \KeywordTok{as.integer} \NormalTok{(nblk)}
    \NormalTok{)}
  \KeywordTok{return} \NormalTok{(h)}
\NormalTok{\}}

\NormalTok{sortBlocks <-}\StringTok{ }\NormalTok{function (y, w, xind, iblks) \{}
  \NormalTok{n <-}\StringTok{ }\KeywordTok{length} \NormalTok{(y)}
  \NormalTok{nblk <-}\StringTok{ }\KeywordTok{max} \NormalTok{(iblks)}
  \NormalTok{h <-}\StringTok{ }\KeywordTok{.C}\NormalTok{(}
    \StringTok{"sortBlocks"}\NormalTok{,}
    \DataTypeTok{y =} \KeywordTok{as.double} \NormalTok{(y),}
    \DataTypeTok{w =} \KeywordTok{as.double} \NormalTok{(w),}
    \DataTypeTok{xind =} \KeywordTok{as.integer} \NormalTok{(xind),}
    \KeywordTok{as.integer}\NormalTok{(iblks),}
    \KeywordTok{as.integer} \NormalTok{(n),}
    \KeywordTok{as.integer} \NormalTok{(nblk)}
  \NormalTok{)}
  \KeywordTok{return} \NormalTok{(h)}
\NormalTok{\}}

\NormalTok{jbkPava <-}\StringTok{ }\NormalTok{function (x, }\DataTypeTok{w =} \KeywordTok{rep}\NormalTok{(}\DecValTok{1}\NormalTok{, }\KeywordTok{length}\NormalTok{(x))) \{}
  \NormalTok{h <-}
\StringTok{    }\KeywordTok{.C}\NormalTok{(}\StringTok{"jbkPava"}\NormalTok{,}
       \DataTypeTok{x =} \KeywordTok{as.double}\NormalTok{(x),}
       \DataTypeTok{w =} \KeywordTok{as.double} \NormalTok{(w),}
       \DataTypeTok{n =} \KeywordTok{as.integer}\NormalTok{(}\KeywordTok{length}\NormalTok{(x)))}
  \KeywordTok{return} \NormalTok{(h)}
\NormalTok{\}}

\NormalTok{monregR <-}
\StringTok{  }\NormalTok{function (x,}
            \NormalTok{y,}
            \DataTypeTok{w =} \KeywordTok{rep} \NormalTok{(}\DecValTok{1}\NormalTok{, }\KeywordTok{length} \NormalTok{(x)),}
            \DataTypeTok{ties =} \DecValTok{1}\NormalTok{) \{}
    \NormalTok{f <-}\StringTok{ }\KeywordTok{sort}\NormalTok{(}\KeywordTok{unique}\NormalTok{(x))}
    \NormalTok{g <-}\StringTok{ }\KeywordTok{lapply}\NormalTok{(f, function (z)}
      \KeywordTok{which}\NormalTok{(x ==}\StringTok{ }\NormalTok{z))}
    \NormalTok{n <-}\StringTok{ }\KeywordTok{length} \NormalTok{(x)}
    \NormalTok{k <-}\StringTok{ }\KeywordTok{length} \NormalTok{(f)}
    \NormalTok{if (ties ==}\StringTok{ }\DecValTok{1}\NormalTok{) \{}
      \NormalTok{h <-}\StringTok{ }\KeywordTok{lapply} \NormalTok{(g, function (x)}
        \NormalTok{y[x])}
      \NormalTok{f <-}\StringTok{ }\KeywordTok{lapply} \NormalTok{(g, function (x)}
        \NormalTok{w[x])}
      \NormalTok{m <-}\StringTok{ }\KeywordTok{rep} \NormalTok{(}\DecValTok{0}\NormalTok{, n)}
      \NormalTok{for (i in }\DecValTok{1}\NormalTok{:k) \{}
        \NormalTok{ii <-}\StringTok{ }\KeywordTok{order} \NormalTok{(h[[i]])}
        \NormalTok{g[[i]] <-}\StringTok{ }\NormalTok{g[[i]][ii]}
        \NormalTok{h[[i]] <-}\StringTok{ }\NormalTok{h[[i]][ii]}
        \NormalTok{f[[i]] <-}\StringTok{ }\NormalTok{f[[i]][ii]}
      \NormalTok{\}}
      \NormalTok{r <-}\StringTok{ }\KeywordTok{jbkPava} \NormalTok{(}\KeywordTok{unlist}\NormalTok{(h), }\KeywordTok{unlist}\NormalTok{(f))$x}
      \NormalTok{s <-}\StringTok{ }\NormalTok{r[}\KeywordTok{order} \NormalTok{(}\KeywordTok{unlist} \NormalTok{(g))]}
    \NormalTok{\}}
    \NormalTok{if (ties ==}\StringTok{ }\DecValTok{2}\NormalTok{) \{}
      \NormalTok{h <-}\StringTok{ }\KeywordTok{lapply} \NormalTok{(g, function (x)}
        \NormalTok{y[x])}
      \NormalTok{f <-}\StringTok{ }\KeywordTok{lapply} \NormalTok{(g, function (x)}
        \NormalTok{w[x])}
      \NormalTok{s <-}\StringTok{ }\KeywordTok{sapply}\NormalTok{(}\DecValTok{1}\NormalTok{:k, function(i) }\KeywordTok{sum} \NormalTok{(h[[i]] *}\StringTok{ }\NormalTok{f[[i]]))}
      \NormalTok{r <-}\StringTok{ }\KeywordTok{sapply} \NormalTok{(f, sum)}
      \NormalTok{m <-}\StringTok{ }\NormalTok{s /}\StringTok{ }\NormalTok{r}
      \NormalTok{r <-}\StringTok{ }\KeywordTok{jbkPava} \NormalTok{(m, r)$x}
      \NormalTok{s <-}\StringTok{ }\KeywordTok{rep} \NormalTok{(}\DecValTok{0}\NormalTok{, n)}
      \NormalTok{for (i in }\DecValTok{1}\NormalTok{:k)}
        \NormalTok{s[g[[i]]] <-}\StringTok{ }\NormalTok{r[i]}
    \NormalTok{\}}
    \NormalTok{if (ties ==}\StringTok{ }\DecValTok{3}\NormalTok{) \{}
      \NormalTok{h <-}\StringTok{ }\KeywordTok{lapply} \NormalTok{(g, function (x)}
        \NormalTok{y[x])}
      \NormalTok{f <-}\StringTok{ }\KeywordTok{lapply} \NormalTok{(g, function (x)}
        \NormalTok{w[x])}
      \NormalTok{s <-}\StringTok{ }\KeywordTok{sapply}\NormalTok{(}\DecValTok{1}\NormalTok{:k, function(i) }\KeywordTok{sum} \NormalTok{(h[[i]] *}\StringTok{ }\NormalTok{f[[i]]))}
      \NormalTok{r <-}\StringTok{ }\KeywordTok{sapply} \NormalTok{(f, sum)}
      \NormalTok{m <-}\StringTok{ }\NormalTok{s /}\StringTok{ }\NormalTok{r}
      \NormalTok{r <-}\StringTok{ }\KeywordTok{jbkPava} \NormalTok{(m, r)$x}
      \NormalTok{s <-}\StringTok{ }\KeywordTok{rep} \NormalTok{(}\DecValTok{0}\NormalTok{, n)}
      \NormalTok{for (i in }\DecValTok{1}\NormalTok{:k)}
        \NormalTok{s[g[[i]]] <-}\StringTok{ }\NormalTok{y[g[[i]]] +}\StringTok{ }\NormalTok{(r[i] -}\StringTok{ }\NormalTok{m[i])}
    \NormalTok{\}}
    \KeywordTok{return} \NormalTok{(s)}
  \NormalTok{\}}

\NormalTok{monregRC <-}\StringTok{ }\NormalTok{function (x,}
                      \NormalTok{y,}
                      \DataTypeTok{w =} \KeywordTok{rep}\NormalTok{(}\DecValTok{1}\NormalTok{, }\KeywordTok{length}\NormalTok{(x)),}
                      \DataTypeTok{ties =} \DecValTok{1}\NormalTok{) \{}
  \NormalTok{a <-}\StringTok{ }\KeywordTok{mySortDouble}\NormalTok{(x, y, w)}
  \NormalTok{b <-}\StringTok{ }\KeywordTok{tieBlock} \NormalTok{(a$x)$iblks}
  \NormalTok{if (ties ==}\StringTok{ }\DecValTok{1}\NormalTok{) \{}
    \NormalTok{c <-}\StringTok{ }\KeywordTok{sortBlocks} \NormalTok{(a$y, a$w, a$xind, b)}
    \NormalTok{d <-}\StringTok{ }\KeywordTok{jbkPava} \NormalTok{(c$y, c$w)$x}
    \NormalTok{e <-}\StringTok{ }\KeywordTok{mySortInteger} \NormalTok{(c$xind)$xind}
    \KeywordTok{return} \NormalTok{(d[e])}
  \NormalTok{\}}
  \NormalTok{if (ties ==}\StringTok{ }\DecValTok{2}\NormalTok{) \{}
    \NormalTok{c <-}\StringTok{ }\KeywordTok{makeBlocks} \NormalTok{(a$y, a$w, b)}
    \NormalTok{d <-}\StringTok{ }\KeywordTok{jbkPava} \NormalTok{(c$xblks, c$wblks)$x}
    \NormalTok{e <-}\StringTok{ }\KeywordTok{mySortInteger} \NormalTok{(a$xind)$xind}
    \KeywordTok{return} \NormalTok{(d[b][e])}
  \NormalTok{\}}
  \NormalTok{if (ties ==}\StringTok{ }\DecValTok{3}\NormalTok{) \{}
    \NormalTok{c <-}\StringTok{ }\KeywordTok{makeBlocks} \NormalTok{(a$y, a$w, b)}
    \NormalTok{d <-}\StringTok{ }\KeywordTok{jbkPava} \NormalTok{(c$xblks, c$wblks)$x}
    \NormalTok{u <-}\StringTok{ }\NormalTok{a$y -}\StringTok{ }\NormalTok{(c$xblks -}\StringTok{ }\NormalTok{d)[b]}
    \NormalTok{e <-}\StringTok{ }\KeywordTok{mySortInteger} \NormalTok{(a$xind)$xind}
    \KeywordTok{return} \NormalTok{(u[e])}
  \NormalTok{\}}
\NormalTok{\}}
\end{Highlighting}
\end{Shaded}

\subsection{C code}\label{c-code}

\begin{Shaded}
\begin{Highlighting}[]

\OtherTok{#include <math.h>}
\OtherTok{#include <stdlib.h>}
\OtherTok{#include <stdbool.h>}

\KeywordTok{struct} \NormalTok{block \{}
    \DataTypeTok{double} \NormalTok{value;}
    \DataTypeTok{double} \NormalTok{weight;}
    \DataTypeTok{int} \NormalTok{size;}
    \DataTypeTok{int} \NormalTok{previous;}
    \DataTypeTok{int} \NormalTok{next;}
\NormalTok{\};}

\KeywordTok{struct} \NormalTok{quadruple \{}
    \DataTypeTok{double} \NormalTok{value;}
    \DataTypeTok{double} \NormalTok{result;}
    \DataTypeTok{double} \NormalTok{weight;}
    \DataTypeTok{int} \NormalTok{index;}
\NormalTok{\};}

\KeywordTok{struct} \NormalTok{triple \{}
    \DataTypeTok{double} \NormalTok{value;}
    \DataTypeTok{double} \NormalTok{weight;}
    \DataTypeTok{int} \NormalTok{index;}
\NormalTok{\};}

\KeywordTok{struct} \NormalTok{pair \{}
    \DataTypeTok{int} \NormalTok{value;}
    \DataTypeTok{int} \NormalTok{index;}
\NormalTok{\};}

\DataTypeTok{int} \NormalTok{myCompDouble (}\DataTypeTok{const} \DataTypeTok{void} \NormalTok{*, }\DataTypeTok{const} \DataTypeTok{void} \NormalTok{*);}
\DataTypeTok{int} \NormalTok{myCompInteger (}\DataTypeTok{const} \DataTypeTok{void} \NormalTok{*, }\DataTypeTok{const} \DataTypeTok{void} \NormalTok{*);}
\DataTypeTok{void} \NormalTok{mySortDouble (}\DataTypeTok{double} \NormalTok{*, }\DataTypeTok{double} \NormalTok{*, }\DataTypeTok{double} \NormalTok{*, }\DataTypeTok{int} \NormalTok{*, }\DataTypeTok{const} \DataTypeTok{int} \NormalTok{*);}
\DataTypeTok{void} \NormalTok{mySortInteger (}\DataTypeTok{int} \NormalTok{*, }\DataTypeTok{int} \NormalTok{*, }\DataTypeTok{const} \DataTypeTok{int} \NormalTok{*);}
\DataTypeTok{void} \NormalTok{mySortInBlock (}\DataTypeTok{double} \NormalTok{*, }\DataTypeTok{double} \NormalTok{*, }\DataTypeTok{int} \NormalTok{*, }\DataTypeTok{int} \NormalTok{*);}
\DataTypeTok{void} \NormalTok{tieBlock (}\DataTypeTok{double} \NormalTok{*, }\DataTypeTok{int} \NormalTok{*, }\DataTypeTok{const} \DataTypeTok{int} \NormalTok{*, }\DataTypeTok{int} \NormalTok{*);}
\DataTypeTok{void} \NormalTok{makeBlocks (}\DataTypeTok{double} \NormalTok{*, }\DataTypeTok{double} \NormalTok{*, }\DataTypeTok{double} \NormalTok{*, }\DataTypeTok{double} \NormalTok{*, }\DataTypeTok{int} \NormalTok{*, }\DataTypeTok{const} \DataTypeTok{int} \NormalTok{*, }\DataTypeTok{const} \DataTypeTok{int} \NormalTok{*);}
\DataTypeTok{void} \NormalTok{sortBlocks (}\DataTypeTok{double} \NormalTok{*, }\DataTypeTok{double} \NormalTok{*, }\DataTypeTok{int} \NormalTok{*, }\DataTypeTok{const} \DataTypeTok{int} \NormalTok{*, }\DataTypeTok{const} \DataTypeTok{int} \NormalTok{*, }\DataTypeTok{const} \DataTypeTok{int} \NormalTok{*);}
\DataTypeTok{void} \NormalTok{jbkPava (}\DataTypeTok{double} \NormalTok{*, }\DataTypeTok{double} \NormalTok{*, }\DataTypeTok{const} \DataTypeTok{int} \NormalTok{*);}

\DataTypeTok{int} \NormalTok{myCompDouble (}\DataTypeTok{const} \DataTypeTok{void} \NormalTok{*px, }\DataTypeTok{const} \DataTypeTok{void} \NormalTok{*py) \{}
    \DataTypeTok{double} \NormalTok{x = ((}\KeywordTok{struct} \NormalTok{quadruple *)px)->value;}
    \DataTypeTok{double} \NormalTok{y = ((}\KeywordTok{struct} \NormalTok{quadruple *)py)->value;}
    \KeywordTok{return} \NormalTok{(}\DataTypeTok{int}\NormalTok{)copysign(}\FloatTok{1.0}\NormalTok{, x - y);}
\NormalTok{\}}

\DataTypeTok{int} \NormalTok{myCompInteger (}\DataTypeTok{const} \DataTypeTok{void} \NormalTok{*px, }\DataTypeTok{const} \DataTypeTok{void} \NormalTok{*py) \{}
    \DataTypeTok{int} \NormalTok{x = ((}\KeywordTok{struct} \NormalTok{pair *)px)->value;}
    \DataTypeTok{int} \NormalTok{y = ((}\KeywordTok{struct} \NormalTok{pair *)py)->value;}
    \KeywordTok{return} \NormalTok{(}\DataTypeTok{int}\NormalTok{)copysign(}\FloatTok{1.0}\NormalTok{, x - y);}
\NormalTok{\}}

\DataTypeTok{void} \NormalTok{mySortInBlock (}\DataTypeTok{double} \NormalTok{*x, }\DataTypeTok{double} \NormalTok{*w, }\DataTypeTok{int} \NormalTok{*xind, }\DataTypeTok{int} \NormalTok{*n) \{}
    \DataTypeTok{int} \NormalTok{nn = *n;}
    \KeywordTok{struct} \NormalTok{triple *xi =}
        \NormalTok{(}\KeywordTok{struct} \NormalTok{triple *) calloc((size_t) nn, (size_t) }\KeywordTok{sizeof}\NormalTok{(}\KeywordTok{struct} \NormalTok{triple));}
    \KeywordTok{for} \NormalTok{(}\DataTypeTok{int} \NormalTok{i = }\DecValTok{0}\NormalTok{; i < nn; i++) \{}
        \NormalTok{xi[i].value = x[i];}
        \NormalTok{xi[i].weight = w[i];}
        \NormalTok{xi[i].index = xind[i];}
    \NormalTok{\}}
    \NormalTok{(}\DataTypeTok{void}\NormalTok{) qsort(xi, (size_t)nn, (size_t)}\KeywordTok{sizeof}\NormalTok{(}\KeywordTok{struct} \NormalTok{triple), myCompDouble);}
    \KeywordTok{for} \NormalTok{(}\DataTypeTok{int} \NormalTok{i = }\DecValTok{0}\NormalTok{; i < nn; i++) \{}
        \NormalTok{x[i] = xi[i].value;}
        \NormalTok{w[i] = xi[i].weight;}
        \NormalTok{xind[i] = xi[i].index;}
    \NormalTok{\}}
    \NormalTok{free(xi);}
    \KeywordTok{return}\NormalTok{;}
\NormalTok{\}}


\DataTypeTok{void} \NormalTok{mySortDouble (}\DataTypeTok{double} \NormalTok{*x, }\DataTypeTok{double} \NormalTok{*y, }\DataTypeTok{double} \NormalTok{*w, }\DataTypeTok{int} \NormalTok{*xind, }\DataTypeTok{const} \DataTypeTok{int} \NormalTok{*n) \{}
    \DataTypeTok{int} \NormalTok{nn = *n;}
    \KeywordTok{struct} \NormalTok{quadruple *xi =}
        \NormalTok{(}\KeywordTok{struct} \NormalTok{quadruple *) calloc((size_t) nn, (size_t) }\KeywordTok{sizeof}\NormalTok{(}\KeywordTok{struct} \NormalTok{quadruple));}
    \KeywordTok{for} \NormalTok{(}\DataTypeTok{int} \NormalTok{i = }\DecValTok{0}\NormalTok{; i < nn; i++) \{}
        \NormalTok{xi[i].value = x[i];}
        \NormalTok{xi[i].result = y[i];}
        \NormalTok{xi[i].weight = w[i];}
        \NormalTok{xi[i].index = i + }\DecValTok{1}\NormalTok{;}
    \NormalTok{\}}
    \NormalTok{(}\DataTypeTok{void}\NormalTok{) qsort(xi, (size_t)nn, (size_t)}\KeywordTok{sizeof}\NormalTok{(}\KeywordTok{struct} \NormalTok{quadruple), myCompDouble);}
    \KeywordTok{for} \NormalTok{(}\DataTypeTok{int} \NormalTok{i = }\DecValTok{0}\NormalTok{; i < nn; i++) \{}
        \NormalTok{x[i] = xi[i].value;}
        \NormalTok{y[i] = xi[i].result;}
        \NormalTok{w[i] = xi[i].weight;}
        \NormalTok{xind[i] = xi[i].index;}
    \NormalTok{\}}
    \NormalTok{free(xi);}
    \KeywordTok{return}\NormalTok{;}
\NormalTok{\}}

\DataTypeTok{void} \NormalTok{mySortInteger (}\DataTypeTok{int} \NormalTok{*x, }\DataTypeTok{int} \NormalTok{*k, }\DataTypeTok{const} \DataTypeTok{int} \NormalTok{*n) \{}
    \DataTypeTok{int} \NormalTok{nn = *n;}
    \KeywordTok{struct} \NormalTok{pair *xi =}
        \NormalTok{(}\KeywordTok{struct} \NormalTok{pair *) calloc((size_t) nn, (size_t) }\KeywordTok{sizeof}\NormalTok{(}\KeywordTok{struct} \NormalTok{pair));}
    \KeywordTok{for} \NormalTok{(}\DataTypeTok{int} \NormalTok{i = }\DecValTok{0}\NormalTok{; i < nn; i++) \{}
        \NormalTok{xi[i].value = x[i];}
        \NormalTok{xi[i].index = i + }\DecValTok{1}\NormalTok{;}
    \NormalTok{\}}
    \NormalTok{(}\DataTypeTok{void}\NormalTok{) qsort(xi, (size_t)nn, (size_t)}\KeywordTok{sizeof}\NormalTok{(}\KeywordTok{struct} \NormalTok{pair), myCompInteger);}
    \KeywordTok{for} \NormalTok{(}\DataTypeTok{int} \NormalTok{i = }\DecValTok{0}\NormalTok{; i < nn; i++) \{}
        \NormalTok{x[i] = xi[i].value;}
        \NormalTok{k[i] = xi[i].index;}
    \NormalTok{\}}
    \NormalTok{free(xi);}
    \KeywordTok{return}\NormalTok{;}
\NormalTok{\}}

\DataTypeTok{void} \NormalTok{tieBlock (}\DataTypeTok{double} \NormalTok{*x, }\DataTypeTok{int} \NormalTok{*iblks, }\DataTypeTok{const} \DataTypeTok{int} \NormalTok{*n, }\DataTypeTok{int} \NormalTok{*nblk) \{}
    \NormalTok{iblks[}\DecValTok{0}\NormalTok{] = }\DecValTok{1}\NormalTok{;}
    \KeywordTok{for} \NormalTok{(}\DataTypeTok{int} \NormalTok{i = }\DecValTok{1}\NormalTok{; i < *n; i++) \{}
        \KeywordTok{if} \NormalTok{(x[i - }\DecValTok{1}\NormalTok{] == x[i]) \{}
            \NormalTok{iblks[i] = iblks[i - }\DecValTok{1}\NormalTok{];}
        \NormalTok{\} }\KeywordTok{else} \NormalTok{\{}
            \NormalTok{iblks[i] = iblks[i - }\DecValTok{1}\NormalTok{] + }\DecValTok{1}\NormalTok{;}
        \NormalTok{\}}
    \NormalTok{\}}
    \NormalTok{*nblk = iblks[*n - }\DecValTok{1}\NormalTok{];}
    \KeywordTok{return}\NormalTok{;}
\NormalTok{\}}

\DataTypeTok{void} \NormalTok{makeBlocks (}\DataTypeTok{double} \NormalTok{*x, }\DataTypeTok{double} \NormalTok{*w, }\DataTypeTok{double} \NormalTok{*xblks, }\DataTypeTok{double} \NormalTok{*wblks, }\DataTypeTok{int} \NormalTok{*iblks, }\DataTypeTok{const} \DataTypeTok{int} \NormalTok{*n, }\DataTypeTok{const} \DataTypeTok{int} \NormalTok{*nblk) \{}
    \KeywordTok{for} \NormalTok{(}\DataTypeTok{int} \NormalTok{i = }\DecValTok{0}\NormalTok{; i < *nblk; i++) \{}
        \NormalTok{xblks[i] = }\FloatTok{0.0}\NormalTok{;}
        \NormalTok{wblks[i] = }\FloatTok{0.0}\NormalTok{;}
    \NormalTok{\}}
    \KeywordTok{for} \NormalTok{(}\DataTypeTok{int} \NormalTok{i = }\DecValTok{0}\NormalTok{; i < *n; i++) \{}
        \NormalTok{xblks [iblks [i] - }\DecValTok{1}\NormalTok{] += w[i] * x[i];}
        \NormalTok{wblks [iblks [i] - }\DecValTok{1}\NormalTok{] += w[i];}
    \NormalTok{\}}
    \KeywordTok{for} \NormalTok{(}\DataTypeTok{int} \NormalTok{i = }\DecValTok{0}\NormalTok{; i < *nblk; i++) \{}
        \NormalTok{xblks[i] = xblks[i] / wblks[i];}
    \NormalTok{\}}
    \KeywordTok{return}\NormalTok{;}
\NormalTok{\}}

\DataTypeTok{void} \NormalTok{sortBlocks (}\DataTypeTok{double} \NormalTok{*y, }\DataTypeTok{double} \NormalTok{*w, }\DataTypeTok{int} \NormalTok{*xind, }\DataTypeTok{const} \DataTypeTok{int} \NormalTok{*iblks, }\DataTypeTok{const} \DataTypeTok{int} \NormalTok{*n, }\DataTypeTok{const} \DataTypeTok{int} \NormalTok{*nblk) \{}
    \DataTypeTok{int} \NormalTok{*nblks = (}\DataTypeTok{int} \NormalTok{*) calloc((size_t) * nblk, }\KeywordTok{sizeof}\NormalTok{(}\DataTypeTok{int}\NormalTok{));}
    \KeywordTok{for} \NormalTok{(}\DataTypeTok{int} \NormalTok{i = }\DecValTok{0}\NormalTok{; i < *n; i++) \{}
        \NormalTok{nblks[iblks[i] - }\DecValTok{1}\NormalTok{]++;}
    \NormalTok{\}}
    \DataTypeTok{int} \NormalTok{k = }\DecValTok{0}\NormalTok{;}
    \KeywordTok{for} \NormalTok{(}\DataTypeTok{int} \NormalTok{i = }\DecValTok{0}\NormalTok{; i < *nblk; i++) \{}
        \DataTypeTok{int} \NormalTok{nn = nblks[i];}
        \NormalTok{(}\DataTypeTok{void}\NormalTok{) mySortInBlock (y + k, w + k, xind + k, &nn);}
        \NormalTok{k += nn;}
    \NormalTok{\}}
    \NormalTok{free (nblks);}
    \KeywordTok{return}\NormalTok{;}
\NormalTok{\}}

\DataTypeTok{void} \NormalTok{jbkPava (}\DataTypeTok{double} \NormalTok{*x, }\DataTypeTok{double} \NormalTok{*w, }\DataTypeTok{const} \DataTypeTok{int} \NormalTok{*n) \{}
    \KeywordTok{struct} \NormalTok{block *blocks = calloc ((size_t) * n, }\KeywordTok{sizeof}\NormalTok{(}\KeywordTok{struct} \NormalTok{block));}
    \KeywordTok{for} \NormalTok{(}\DataTypeTok{int} \NormalTok{i = }\DecValTok{0}\NormalTok{; i < *n; i++) \{}
        \NormalTok{blocks[i].value = x[i];}
        \NormalTok{blocks[i].weight = w[i];}
        \NormalTok{blocks[i].size = }\DecValTok{1}\NormalTok{;}
        \NormalTok{blocks[i].previous = i - }\DecValTok{1}\NormalTok{; }\CommentTok{// index first element previous block}
        \NormalTok{blocks[i].next = i + }\DecValTok{1}\NormalTok{;     }\CommentTok{// index first element next block}
    \NormalTok{\}}
    \DataTypeTok{int} \NormalTok{active = }\DecValTok{0}\NormalTok{;}
    \KeywordTok{do} \NormalTok{\{}
        \NormalTok{bool upsatisfied = false;}
        \DataTypeTok{int} \NormalTok{next = blocks[active].next;}
        \KeywordTok{if} \NormalTok{(next == *n) upsatisfied = true;}
        \KeywordTok{else} \KeywordTok{if} \NormalTok{(blocks[next].value > blocks[active].value) upsatisfied = true;}
        \KeywordTok{if} \NormalTok{(!upsatisfied) \{}
            \DataTypeTok{double} \NormalTok{ww = blocks[active].weight + blocks[next].weight;}
            \DataTypeTok{int} \NormalTok{nextnext = blocks[next].next;}
            \NormalTok{blocks[active].value = (blocks[active].weight * blocks[active].value + blocks[next].weight * blocks[next].value) / ww;}
            \NormalTok{blocks[active].weight = ww;}
            \NormalTok{blocks[active].size += blocks[next].size;}
            \NormalTok{blocks[active].next = nextnext;}
            \KeywordTok{if} \NormalTok{(nextnext < *n)}
                \NormalTok{blocks[nextnext].previous = active;}
            \NormalTok{blocks[next].size = }\DecValTok{0}\NormalTok{;}
        \NormalTok{\}}
        \NormalTok{bool downsatisfied = false;}
        \DataTypeTok{int} \NormalTok{previous = blocks[active].previous;}
        \KeywordTok{if} \NormalTok{(previous == -}\DecValTok{1}\NormalTok{) downsatisfied = true;}
        \KeywordTok{else} \KeywordTok{if} \NormalTok{(blocks[previous].value < blocks[active].value) downsatisfied = true;}
        \KeywordTok{if} \NormalTok{(!downsatisfied) \{}
            \DataTypeTok{double} \NormalTok{ww = blocks[active].weight + blocks[previous].weight;}
            \DataTypeTok{int} \NormalTok{previousprevious = blocks[previous].previous;}
            \NormalTok{blocks[active].value = (blocks[active].weight * blocks[active].value + blocks[previous].weight * blocks[previous].value) / ww;}
            \NormalTok{blocks[active].weight = ww;}
            \NormalTok{blocks[active].size += blocks[previous].size;}
            \NormalTok{blocks[active].previous = previousprevious;}
            \KeywordTok{if} \NormalTok{(previousprevious > -}\DecValTok{1}\NormalTok{)}
                \NormalTok{blocks[previousprevious].next = active;}
            \NormalTok{blocks[previous].size = }\DecValTok{0}\NormalTok{;}
        \NormalTok{\}}
        \KeywordTok{if} \NormalTok{((blocks[active].next == *n) && downsatisfied) }\KeywordTok{break}\NormalTok{;}
        \KeywordTok{if} \NormalTok{(upsatisfied && downsatisfied) active = next;}
    \NormalTok{\} }\KeywordTok{while} \NormalTok{(true);}
    \DataTypeTok{int} \NormalTok{k = }\DecValTok{0}\NormalTok{;}
    \KeywordTok{for} \NormalTok{(}\DataTypeTok{int} \NormalTok{i = }\DecValTok{0}\NormalTok{; i < *n; i++) \{}
        \DataTypeTok{int} \NormalTok{blksize = blocks[i].size;}
        \KeywordTok{if} \NormalTok{(blksize > }\FloatTok{0.0}\NormalTok{) \{}
            \KeywordTok{for} \NormalTok{(}\DataTypeTok{int} \NormalTok{j = }\DecValTok{0}\NormalTok{; j < blksize; j++) \{}
                \NormalTok{x[k] = blocks[i].value;}
                \NormalTok{k++;}
            \NormalTok{\}}
        \NormalTok{\}}
    \NormalTok{\}}
    \NormalTok{free (blocks);}
\NormalTok{\}}
\end{Highlighting}
\end{Shaded}

\section*{References}\label{references}
\addcontentsline{toc}{section}{References}

\hypertarget{refs}{}
\hypertarget{ref-cran_80}{}
Cran, G.W. 1980. ``Algorithm AS 149: Amalgamation of Means in the Case
of Simple Ordering.'' \emph{Journal of the Royal Statistical Society,
Series C (Applied Statistics)}, 209--11.

\hypertarget{ref-deleeuw_A_77}{}
De Leeuw, J. 1977. ``Correctness of Kruskal's Algorithms for Monotone
Regression with Ties.'' \emph{Psychometrika} 42: 141--44.
\url{http://www.stat.ucla.edu/~deleeuw/janspubs/1977/articles/deleeuw_A_77.pdf}.

\hypertarget{ref-deleeuw_E_16d}{}
---------. 2016. ``Exceedingly Simple Isotone Regression with Ties.''
doi:\href{https://doi.org/10.13140/RG.2.1.3698.2801}{10.13140/RG.2.1.3698.2801}.

\hypertarget{ref-deleeuw_E_17g}{}
---------. 2017a. ``Exceedingly Simple Monotone Regression.''
\url{http://rpubs.com/deleeuw/262796}.

\hypertarget{ref-deleeuw_E_17f}{}
---------. 2017b. ``Exceedingly Simple Sorting with Indices.''
\url{http://rpubs.com/deleeuw/261074}.

\hypertarget{ref-kruskal_64b}{}
Kruskal, J.B. 1964. ``Nonmetric Multidimensional Scaling: a Numerical
Method.'' \emph{Psychometrika} 29: 115--29.


\end{document}
